\clearpage
\begingroup
\tcbset{myboxstyle/.append style={breakable=false}}
\everymath{\textstyle}
\subsection{Субаддитивность и супераддитивность}

Пусть $\mathcal X$ --- система подмножеств некоторого множества,
а $m\colon\mathcal X\to[0;+\infty)$ --- неотрицательная функция.

\begin{definition}[Субаддитивность]
    Функция $m$ называется \textbf{субаддитивной}, если для любых
    $A,A_1,\ldots,A_n\in\mathcal X$
    \[
        A\subseteq\bigcup_{i=1}^{n}A_i
        \quad\Longrightarrow\quad
        m(A)\leq\sum_{i=1}^{n}m(A_i).
    \]
\end{definition}

\begin{definition}[Супераддитивность]
    Функция $m$ называется \textbf{супераддитивной}, если для любых
    $B,B_1,\ldots,B_r\in\mathcal X$, где $B_j$ попарно не пересекаются,
    \[
        \bigsqcup_{j=1}^{r}B_j\subseteq B
        \quad\Longrightarrow\quad
        \sum_{j=1}^{r}m(B_j)\leq m(B).
    \]
\end{definition}

\begin{theorem}
    Любая мера на полукольце $S$ субаддитивна и супераддитивна.
\end{theorem}

\begin{Proof}[Субаддитивность]
    Пусть $A\subseteq\bigcup_{i=1}^{n}A_i$, где $A,A_i\in S$.
    По лемме о замельчении существует конечное семейство попарно
    непересекающихся множеств $C_j\in S$, для которого
    \[
        A=\bigsqcup_{j\in I}C_j,\qquad
        A_i=\bigsqcup_{j\in I_i}C_j.
    \]
    Каждое $C_j$, входящее в разбиение $A$, входит хотя бы в одно разбиение
    $A_i$. Поэтому по конечной аддитивности и неотрицательности
    \[
        m(A)=\sum_{j\in I}m(C_j)
        \leq\sum_{i=1}^{n}\sum_{j\in I_i}m(C_j)
        =\sum_{i=1}^{n}m(A_i).
    \]
\end{Proof}

\begin{Proof}[Супераддитивность]
    По свойствам полукольца найдутся попарно непересекающиеся $D_k\in S$:
    \[
        B\setminus\bigsqcup_{j=1}^{r}B_j
        =\bigsqcup_{k=1}^{\ell}D_k,\qquad D_k\in S.
    \]
    По аддитивности $m(B)=\sum_{j=1}^{r}m(B_j)+\sum_{k=1}^{\ell}m(D_k)
    \geq\sum_{j=1}^{r}m(B_j)$.
\end{Proof}

\newpage
\begin{theorem}[Обратное утверждение]
    Если неотрицательная функция $m$ на полукольце $S$ субаддитивна
    и супераддитивна, то она конечно аддитивна, то есть является мерой.
\end{theorem}

\begin{Proof}
    При $A=\bigsqcup_{i=1}^{n}A_i$, $A,A_i\in S$, два свойства дают
    \[
        m(A)\leq\sum_{i=1}^{n}m(A_i)\leq m(A).
    \]
    Следовательно, обе оценки обращаются в равенство.
\end{Proof}

\subsection{Счётные аналоги}

\begin{definition}[$\sigma$-субаддитивность]
    Функция $m$ называется \textbf{$\sigma$-субаддитивной}, если для любых
    $A,A_1,A_2,\ldots\in\mathcal X$
    \[
        A\subseteq\bigcup_{i=1}^{\infty}A_i
        \quad\Longrightarrow\quad
        m(A)\leq\sum_{i=1}^{\infty}m(A_i).
    \]
\end{definition}

\begin{definition}[$\sigma$-супераддитивность]
    Функция $m$ называется \textbf{$\sigma$-супераддитивной}, если для любых
    $B,B_1,B_2,\ldots\in\mathcal X$, где $B_j$ попарно не пересекаются,
    \[
        \bigsqcup_{j=1}^{\infty}B_j\subseteq B
        \quad\Longrightarrow\quad
        \sum_{j=1}^{\infty}m(B_j)\leq m(B).
    \]
\end{definition}

\begin{theorem}
    Любая мера на полукольце $S$ является $\sigma$-супераддитивной.
\end{theorem}

\begin{Proof}
    Если $\bigsqcup_{j=1}^{\infty}B_j\subseteq B$, то для каждого $n\in\N$
    \[
        \bigsqcup_{j=1}^{n}B_j\subseteq B,
        \qquad\sum_{j=1}^{n}m(B_j)\leq m(B).
    \]
    Переходя к пределу при $n\to\infty$, получаем требуемое неравенство.
\end{Proof}

\newpage
\subsection{Критерий счётной аддитивности}

\begin{theorem}
    Мера $m$ на полукольце $S$ является $\sigma$-аддитивной
    тогда и только тогда, когда она $\sigma$-субаддитивна.
\end{theorem}

\begin{Proof}[Достаточность]
    Пусть $m$ является $\sigma$-субаддитивной. Если
    $A=\bigsqcup_{i=1}^{\infty}A_i$, где $A,A_i\in S$, то
    по $\sigma$-субаддитивности и $\sigma$-супераддитивности
    \[
        m(A)\leq\sum_{i=1}^{\infty}m(A_i)\leq m(A).
    \]
    Значит, $m(A)=\sum_{i=1}^{\infty}m(A_i)$.
\end{Proof}

\begin{Proof}[Необходимость]
    Пусть $m$ является $\sigma$-аддитивной и
    $A\subseteq\bigcup_{i=1}^{\infty}A_i$, где $A,A_i\in S$.
    Положим $A_i'=A\cap A_i\in S$. Тогда
    \[
        A=\bigcup_{i=1}^{\infty}A_i',\qquad m(A_i')\leq m(A_i).
    \]
    Исключим из каждого множества уже учтённые части. Для $i\geq2$
    представим
    \[
        A_i'\setminus\bigcup_{k=1}^{i-1}A_k'
        =\bigsqcup_{j=1}^{r_i}B_{ij},\qquad B_{ij}\in S.
    \]
    Все полученные множества попарно не пересекаются, и
    \[
        A=A_1'\sqcup\bigsqcup_{i=2}^{\infty}
                       \bigsqcup_{j=1}^{r_i}B_{ij}.
    \]
    По конечной супераддитивности
    $\sum_{j=1}^{r_i}m(B_{ij})\leq m(A_i')$.
    Поэтому по $\sigma$-аддитивности
    \[
        \begin{aligned}
            m(A)
            &=m(A_1')+\sum_{i=2}^{\infty}\sum_{j=1}^{r_i}m(B_{ij})\\
            &\leq\sum_{i=1}^{\infty}m(A_i')
             \leq\sum_{i=1}^{\infty}m(A_i).
        \end{aligned}
    \]
\end{Proof}

\newpage
\subsection{Счётная аддитивность меры Лебега--Стилтьеса}

Пусть $S$ --- полукольцо промежутков, введённое ранее, а
$F\colon\R\to\R$ ограничена, не убывает и непрерывна справа.
Мера $m$ задаётся равенством $m((a;b])=F(b)-F(a)$
с соответствующими предельными формулами для неограниченных промежутков.

\begin{theorem}
    Мера Лебега--Стилтьеса $m$ на полукольце $S$ является
    $\sigma$-аддитивной.
\end{theorem}

\begin{Proof}[Счётная субаддитивность на ограниченном промежутке]
    По доказанному критерию достаточно проверить $\sigma$-субаддитивность.
    Пусть
    \[
        (a;b]\subseteq\bigcup_{i=1}^{\infty}(a_i;b_i].
    \]
    Зафиксируем $\eps>0$. По непрерывности $F$ справа выберем
    $a'\in(a;b)$ и $b_i'>b_i$ так, чтобы
    \[
        F(a')-F(a)<\frac\eps2,\qquad
        F(b_i')-F(b_i)<\frac{\eps}{2^{i+1}}.
    \]
    Открытые интервалы $(a_i;b_i')$ покрывают отрезок $[a';b]$.
    По лемме Гейне--Бореля существует конечное подпокрытие
    с множеством индексов $I\subset\N$:
    \[
        [a';b]\subseteq\bigcup_{i\in I}(a_i;b_i').
    \]
    Следовательно, $(a';b]\subseteq\bigcup_{i\in I}(a_i;b_i']$.
    По конечной субаддитивности
    \[
        \begin{aligned}
            F(b)-F(a')
            &\leq\sum_{i\in I}\bigl(F(b_i')-F(a_i)\bigr)\\
            &\leq\sum_{i=1}^{\infty}
                 \left(F(b_i)-F(a_i)+\frac{\eps}{2^{i+1}}\right)\\
            &=\sum_{i=1}^{\infty}m((a_i;b_i])+\frac\eps2.
        \end{aligned}
    \]
    Прибавляя $F(a')-F(a)<\eps/2$, получаем
    \[
        m((a;b])\leq\sum_{i=1}^{\infty}m((a_i;b_i])+\eps.
    \]
    Устремляя $\eps$ к нулю, получаем требуемое неравенство.
\end{Proof}

\begin{remark}[Неограниченные промежутки]
    Для покрытия лучами используется то же рассуждение с открытыми лучами.
    Если покрываемый промежуток неограничен, сначала рассматривают его
    пересечение с $(-M;M]$, затем переходят к пределу при $M\to\infty$
    по формулам через $F(-\infty)$ и $F(+\infty)$.
\end{remark}
\endgroup
